\section{Bug Bounty 与 Safe Harbor：奖励不能替代授权边界}

VDP 与 bug bounty 经常一起出现，却解决不同问题。\term{bug bounty program}（漏洞赏金计划）是在既定规则下对符合条件的有效报告提供金钱、积分或认可的可选计划；VDP 即使没有奖金也应存在。Bounty 帮助竞争研究注意力、设定 severity rewards 和运营 triage，但不能通过事后付款重写已经发生的有害行为。

\subsection{VDP 与 bounty 的边界}

一个组织可以有 VDP 而没有 bounty，也可以通过第三方平台管理 bounty。两者共享资产范围、报告渠道与协调机制，但 bounty 还需要 reward table、duplicate rule、eligibility、tax/payment 与 appeal。安全领导者必须防止团队因为“这是 bounty 报告”而忽略 data access、extortion signal 或其他 incident indicators。

\lecturefigure{04-vdp-vs-bounty.png}{VDP 提供安全报告渠道，bug bounty 在此基础上增加可选奖励机制。}{CISA/DOJ VDP guidance；本地字幕 12:30--15:20；概念重绘。}

\begin{knowledgebox}{读图：先问授权，再问奖励}
VDP 的核心是 safe channel、scope、response promise 和 coordinated disclosure；bounty 的额外层是 optional payment、severity table、eligibility 与 duplicate。二者都不覆盖 extortion、破坏、超出必要范围的数据访问或拒绝停止测试。评估报告时先判断行为是否符合预先公布的边界和 good faith，再判断漏洞价值和奖励，不要让支付流程决定法律或事件分类。
\end{knowledgebox}

\begin{warningbox}{Payment 不是 post-hoc authorization}
组织支付奖金、签署 NDA 或继续沟通，可能是为了控制风险和促成修复；这些动作不会自动证明此前行为获得授权，也不能消除保存证据、通知受影响方或咨询 counsel 的需要。任何“付钱就把 incident 变成 bounty”的捷径都会制造治理盲点。
\end{warningbox}

\subsection{Safe harbor、scope 与 good faith}

\term{safe harbor}（安全港）在 VDP 语境中是组织对符合政策、善意研究行为的承诺，例如不主张某些民事权利、支持研究者说明其善意，或在适用范围内不建议起诉。它不是普遍法律豁免，也不能约束所有第三方。\term{scope}（范围）明确可测试的域名、应用、API、版本与排除资产；\term{good faith}（善意）通常要求只做验证所需的最小访问、避免隐私和服务影响、停止并及时报告。

\lecturefigure{05-safe-harbor.png}{Safe harbor 只有与清晰 scope、善意要求、禁止行为和升级路径结合才有效。}{CISA VDP template；DOJ VDP framework；概念重绘。}

\begin{knowledgebox}{读图：政策必须同时保护研究者与用户}
In-scope 列出允许测试的资产和方法；good faith 要求最小化访问、报告和保护敏感数据；prohibited 明确 disruption、persistence、data use 或 extortion 等边界；escalation 提供 safety contact、legal contact 与争议路径。只有鼓励语没有具体 scope，会让研究者承担不可预测风险；只有禁止清单没有响应承诺，则不会形成信任。
\end{knowledgebox}

\begin{importantbox}{政策设计检查表}
资产清单必须可维护，第三方服务必须标出所有权边界，测试速率与自动化限制应解释原因，敏感数据处理应规定停止、最小保留和删除，争议应有人工 escalation。政策更新还应保留版本与生效时间，避免组织事后用新规则评价旧行为。
\end{importantbox}

\subsection{本章小结}

VDP 管渠道，bounty 管激励，safe harbor 管符合边界的组织承诺。三者都不能替代 incident classification，也不能把支付、NDA 或善意假设当成技术事实。

